\scanpage{44}
\mspage{18}
% archive foliation: 46
b') Cas $\mathcal X\simeq\mathbb P^1_{\mathbb C}$, $\sigma$ une inversion --- laquelle définit une structure de droite projective réelle --- on trouve \deleted{le disque} pour $X$ un \uline{disque}. Pour obtenir une métrique distinguée sur $\mathcal X$ invariante par $\sigma$, il faut se donner un pt du disque ouvert $X\setminus\partial X$

\figure{\begin{tikzpicture}\draw (0,0)--(5,.2);\end{tikzpicture}}

En résumé, tte surface conforme à bord 0-connexe admet une unique métrique Riemannienne distinguée, sauf dans les quatre cas suivants

a) a') Cas $X\simeq\mathbb C$, ou $X\simeq\unsure{\overline{\mathbb H}|\mathbb H}=\{z\in\mathbb C\mid\mathrm{Im}\,z\geq0\}$ [droite affine complexe, demi-plan à bord défini par une droite affine réelle], on a alors une classe de métriques mod constante, qui forment une orbite sous $\mathrm{Aut}\,X$.

b) b') $X\simeq\mathbb P^1_{\mathbb C}$ ou $X\simeq\text{disque}$ [droite projective complexe ou \deleted{\unsure{demidisque fermé}} défini par droite projective réelle] --- il y a alors un ens. de métriques distinguées dépendant de 3 paramètres réels resp. de 2 paramètres réels, également orbite sous $\mathrm{Aut}\,X$.
